\scanpage{387}\mspage{302 bis}

\deleted{\unsure{devient}} \deleted{\unsure{est}} les $J_r$ stables dans $I$ pour ces limites,
donc il en est en particulier ainsi pour les $\varinjlim$
de type $I_{\alpha_c}$. On peut donc appliquer le résultat
4.13.10.2 pour conclure que si les $J_r$ sont
cofinaux \deleted{\unsure{ds}} $I$ (qui est grd devant $\pi$, $\pi$ étant utile
pour $M$), leur $J=\underline S_{/x}$ l'est également. Cela implique
alors, comme on a vu, que $\underline S$ est $\pi$-accessible
dans $M$\unsure{!|.}

\uline{4.13.10.3.} On peut reformuler la reformulation
précédente, en utilisant la terminologie introduite
plus bas (4.13.23) :

Soit $c$ un cardinal, \deleted{\unsure{non limite}} $\alpha_c$ le plus petit
ordinal $\alpha$ tel que $\operatorname{card}\alpha=c$, $\pi$ un cardinal $\geq c$.

On suppose

$(*)$ \deleted{\unsure{que}} $\pi$ est \inserted{utile} pour $M$, et $M$ est stable \unsure{par|pour} $\varinjlim$ de
type $I_{\alpha_c}$. (\inserted{\deleted{\unsure{utile}} \unsure{admet un card.}}c'est le cas p. ex. si \inserted{$c$ est \uline{non limite} et} $c>\pi_0$, $\pi_0$ étant un
cardinal \deleted{\unsure{utile}} tel que $M$ soit grd $\unsure{\mathrm{dt}|\mathrm{dvt}}$ $\pi_0$, p. ex. un
cardinal \uline{utile} pour $M^{(*)}$.)
% slanted
\margin{$^{(*)}$ Car alors $I_{\alpha_c}$ est grd devant $\pi_0$}

Alors l'ensemble des sous-catégories $\underline S$ de $M$ qui
sont $\alpha_c$-$\pi$-accessibles dans $M$ (i.e. qui sont
$\pi$-accessibles et stables dans $M$ par $\varinjlim$ de type $I_{\alpha_c}$)
est stable par intersections de cardinal $\leq c$. Plus
généralement, si $\unsure{\mathscr L|\underline{\mathscr L}}$ est un ens. de catégories
telles que $1^\circ)$ $I\in\mathscr L\Longrightarrow\operatorname{Card} I\leq\pi$ et $2^\circ)$ $I_{\alpha_c}\in\mathscr L$,
alors l'ens. des sous-catégories de $M$ qui sont
$\mathscr L$-$\pi$-accessibles dans $M$, \unsure{est} stable par \unsure{l'intersection|intersections}
\unsure{de|en} cardinal $\leq c$.
% archive sheet number: 386
